\section{Colombia's payments setting}
Bre-B entered full operation on 6 October 2025, following an implementation period in which private instant-payment initiatives and participating providers were already active \citep{banrep2026}. The central bank describes Bre-B as a service for domestic payments and transfers through participating financial institutions, not an application or an international remittance service \citep{banrepfaq}. Our survey interviews began after full operation, but the survey does not identify which rail processed a reported transaction. The calendar does not establish exposure or create a pre/post comparison.

Interoperability addresses whether providers can exchange a payment. It does not establish that a customer can enroll under her own name, navigate an interface without help, understand a charge or get a problem resolved. These are separate questions about implementation. Colombia's earlier PEP and later PPT were distinct regularization episodes; neither the recorded survey responses nor the authors' account of an onboarding difficulty verifies a participant's legal status or a provider's compliance.

The modern payment environment predates Bre-B. The framework for specialized electronic deposits and payments firms dates to 2014 \citep{law1735}. The 2021 protection statute creates a distinct temporary permit for eligible Venezuelan migrants \citep{decree216}. A legally recognized credential, a provider's technical ability to validate it, and a customer's continued use of her account are separate questions. The official history of Bre-B distinguishes earlier provider initiatives from full operation. Our recall windows apply to the survey questions, not to the launch date; they do not identify when an individual first used the rail.
